\chapter{Ramsey 理论}
\label{ch:250421}

\section{Ramsey Number}

如果 $n$ 阶图的 Y / B 边着色一定有 Y 色的 $K_y$ 或 B 色的 $K_b$ ，我们记作 $n\to(b,y)$. 
显然如果 $n\to(b,y)$，那么 $n+1\to(b,y)$.

首先我们有 $r(x,2)=r(2,x)=x$.

\paragraph*{定理1: $6\to(3,3)$}

随便取一个点 $v$，这个点度数是 5. 在这个点有关的 5 条边中一定有 3 条同色的边 $va,vb,vc$.

如果这 3 条边都是 Y 色的，那么要么 $a,b,c$ 构成 B 色的 $K_3$，要么 $a,b,c$ 两两之间存在 Y 色的边，
这条 Y 色边和 $v$ 构成 Y 色的 $K_3$.  B 色同理.

\paragraph*{定理1.5: $9\to(3,4)$}

注意每个点度数是 $8$. 如果存在 $v$ 有 $4$ 条 b 边, 那对应的 4 个邻接点要么有 B 色 $K_2$, 要么有 Y 色 $K_4$；

如果存在 $v$ 有 6 条 y 边, 那对应的 $6$ 个邻接点要么有 B 色 $K_3$, 要么有 Y 色 $K_3$.

如果以上两种都不成立，那么所有点都是 3 条 b 边和 5 条 y 边. 
我们考虑 B 色这个子图，发现违背 shaking hands lemma（3*9 是奇数）.

\paragraph*{定理2: $18\to(4,4)$}

注意每个点都有 17 条边，其中一定有 9 条边同色.

如果这 9 个边都是 B 色，那这 9 个点要么有 B 色 $K_3$, 要么有 Y 色 $K_4$;

如果这 9 个边都是 Y 色，那这 9 个点要么有 Y 色 $K_3$, 要么有 B 色 $K_4$.

\section{定理3 \& 定理4}

\begin{theorem}
    对于任意正整数 $y,b$，存在一个正整数 $N$ 满足：对于任意 $K_N$ 的 Y / B 边着色，必有 Y 色的 $K_y$ 或 B 色 $K_b$.

    满足条件的最小整数 $N$ 称为 Ramsey 数，记作 $r(y,b)$.
\end{theorem}

要证 $r(y,b)\leq N$ 即证任意 $K_N$  Y / B 边着色，必有 Y 色的 $K_y$ 或 B 色 $K_b$.

要证 $r(y,b)> M$ 即证存在 $K_M$  Y / B 边着色，没有 Y 色的 $K_y$ ，也没有 B 色 $K_b$ （或存在 $M$ 阶图 $G$ ，$\omega(G)<y$ 且 $\alpha(G)<b$）.

\begin{proof}
    归纳法。首先有 $r(2,x)=r(x,2)=x$. 对 $y+b$ 归纳，如果我们有 $r(y-1,b)+r(y,b-1)$ 个顶点，那么每个点度数 $r(y-1,b)+r(y,b-1)-1$.

    因此对任意一个点，要么其 Y 色边至少 $r(y-1,b)$ 条，要么其 B 色边至少 $r(y,b-1)$ 条. 
    （如果 Y 色边小于 $r(y-1,b)$，B 色边小于 $r(y,b-1)$，加起来都不到度数，矛盾）.

    所以 $r(y-1,b)+r(y,b-1)\to(y,b)$ 也就是说 $r(y,b)\leq r(y-1,b)+r(y,b-1)$.
\end{proof}

\begin{lemma}
Let \(q,\gamma_2,\delta=2^d,\beta\) be positive integers such that \(q=(2\gamma_2\delta+1)\), and \(\beta<\gamma_2\). 
Let \(\widetilde{\mathcal{U}}\)
be the following distribution with support
\(
W_1=\{0,1,\ldots,\delta-1\}:
\)
\begin{equation}
\widetilde{\mathcal{U}}(a_1,a_0')=??
% \begin{cases}
% \dfrac{2\beta+1}{2(\beta+1)\delta+1},
%     & \text{if } a_1\in W_1\setminus\{0\},\\[6pt]
% \dfrac{2\beta+2}{2(\beta+1)\delta+1},
%     & \text{if } a_1=0.
% \end{cases}
\end{equation}

For an arbitrary integer \(c_e\) such that
\[
|c_e|\leq \beta<\gamma_2,
\]
\(\widetilde{\mathcal{U}}\) is perfectly indistinguishable from the
following distribution \(\mathcal{U}_{c_e}\): sample uniform \(w'\in\mathbb{Z}_q\)
conditioned on
\[
|\mathrm{Lowbits}(w')|\geq \gamma_2-\beta,
\]
then output \(\left(\mathrm{HighBits}(w'+c_e),\mathrm{LowBits}(w')\right)\).
\end{lemma}


\section{证明 $r(4,4)=18$}

即证，存在 $K_{17}$ 的 b/y 染色不存在同色 $K_4$. 每个点连16条边，如果9b+7y，9个点二染色必有3b or 4y, 都会形成同色 K4. 因此每个点都是 8b+8y

正17边形排在圆上，距离1248->b, 3567->y. 先考虑同色三角形，8+8=1，1+1=2，2+2=4，4+4=8，3+3=6，5+5=7，6+6=5，7+7=3，因此等腰。

\section{证明 $r(3,3,3)=17$}

K16， 每个点15条边，如果出现一个点连6条同色边，假设同r色，由 $r(3,3)=6$, 1.6点的边有r，有同色K3；2.6点的边无r，因此b/y染色，必有同b/y色K3.

法1

\notefigure{notes/250421-01.jpg}

法2 考虑 $V=\{0,1\}^5$

\section{多染色}

$r(a_1,a_2,a_3,a_4,a_5)\leq 5^{(a_1-1)+(a_2-1)+(a_3-1)+(a_4-1)+(a_5-1)+1}:=n$

每个点来说，存在一种颜色b，这个点连接的b色边有大于n/5条，因此去掉所有剩下的边和点。这个游戏最多可以进行 $(a_1-1)+(a_2-1)+(a_3-1)+(a_4-1)+(a_5-1)+1$ 轮，剩下同样数量的点，其中每个点到其他点的边都是同色的. 因此存在a1个c1，或a2个c2，或..., 这些存在每个都会形成一个同色 $K_{c_k}$.

\section{证明 $r(k,k)=\Omega(k^2)$}

$(k-1)\times K_{k-1}$ 内部染y，国际边染b

(1971, 构造性证明) $r(k,k)=\Omega(k^3)$

\section{(Erdős, 1947) $r(k,k)>2^{k/2},k\geq3$}

$n\leq 2^{k/2}$ yb染色 每k个点形成一个点集 $T$, 坏集合的概率为 $2^{1-{k\choose 2}}$, 则随机染色，坏集合个数的期望 ${n\choose k}2^{1-{k\choose 2}}$

列表：row—染色方案；col—k元子集，如果第i个k元子集在第j个染色下同色，打✔. 按col统计，每个k元子集有$2^{1+{n\choose 2}-{k\choose 2}}$个✔，共计${n\choose k}2^{1+{n\choose 2}-{k\choose 2}}$. 平均下来，每个染色方案平均有${n\choose k}2^{1-{k\choose 2}}$个✔.

如果${n\choose k}2^{1-{k\choose 2}}< 1$那么存在染色，没有✔.

${n\choose k}2^{1-{k\choose 2}}<\frac{2^{k^2/2+1-{k\choose2}}}{k!}=\frac{2^{k/2+1}}{k!}\leq_{k=3}0.95$

Another example

\notefigure{notes/250421-02.png}

\section{Kolmogrov complexity \&\& 回顾 probabilistic method}

$K(x)\leq |x|+c$ A Turing Machine that reads the input $x$ and output $x$.

如何给出一个足够 random 的字符串？如果可以比较快抓出来某个字符串，那么这种方法就是不够 ramdom 的.

一个精心构造的证明(即，一种染色)同样不够random，以至于大多数人构造出 $X=\sum X_T$ 都非常大.

\section{进一步}

$\sqrt{2}<r(k,k)^{1/k}<4$, thus $\sqrt{2}\leq \liminf_{k\to+\infty}r(k,k)^{1/k}\leq\limsup_{k\to+\infty}r(k,k)^{1/k}\leq4$

Prob0. optimize the bound? (or refine the $\leq$ ?)

Prob1. Does there exist limit?

Prob2. lim = ?

\section{Ramsey Problem on Integers}

\begin{quotation}
\textbf{Thm.} (Schur, also known as \textbf{broken chalk theorem}, 2025.04.21 15:28:03) $\forall c\;\exists N \;\forall \phi:[N]\to[c],\;\exists x+y=z$ with the same color. The minimum $N=S(c)$.
\end{quotation}

$S(1)=2,S(2)=5$

假设 $\phi$ 是任意染色，构造完全图 $K_{N+1}$ 排列在一条线上，编号1-n+1，$\sigma(i,j)=\phi(|i-j|)$. 由 Ramsey，$N\geq r_c(3)-1$, 存在同色三角形, 即 $x+y=z$.  $\blacksquare$

variation：要求 $x\not=y$

方法1：每种颜色分裂成两个，$r_{2c}(3)$

方法2：$r_c(4)$，4点同色，$a+b+c=d$. 若 $a=b$, 选 $x=a,y=b+c$.

\begin{quotation}
\textbf{Ext1.} $K_N$ b/y染色，已知 Y 色边占总的 90\%，如果 N 足够大, 是否任意满足条件的染色都存在 Y 色 $K_{100}$？

实际上，由 Turan，要使得存在 $K_{100}$, 至少 y 边得占 $98/99$
\end{quotation}

\begin{quotation}
\textbf{Ext2.} 只关注y色，如果 $N$ 足够大，从1-N挑10\%，是否不能避免存在等差数列？or 存在3长度等差数列 (Erdős-Turan conjecture)
\end{quotation}


\section{Ramsey equation}

\begin{quotation}
\textbf{Def.} 整数方程 $E$ 是 Ramsey 的，当且仅当 $\forall c\;\exists N\;\forall \phi:[N]\to[c]\;\exists$ 一个 $E$ 的同色解
\end{quotation}

Trivially, $x+2y=3z$ 等等是 Ramsey 的；$x=py$ 不是.

首先 $x_1+\dots+x_{k-1}=x_k$ 一定是 Ramsey 的，因为对任意 $c$ 考虑 $r(k,\dots,k)$, c个k，相当于这么大的图, c染色必有同色 $K_k$, 而同色 $K_k$ 自然包含了一组解（每条边的差值对应了某个 $x_i$）。

\subsection{$x+ty=z$ 是 Ramsey 方程}

从 $x+6y=z$ 开始, 我们想到了 van der Waerden 定理：

\begin{quotation}
\textbf{Thm.} (van der Waerden) $\forall c,k\;\exists N\;\forall\phi:[N]\to[c]\;\exists$ a monochromatic AP of length $k$ in $[n]$. The least $N$ is $W(c,k)$.
\end{quotation}

$\forall c\;\exists N\;\forall\phi:[N]\to[c]\;\exists x,d$, 所有 $x,x+d,\dots，x+6kd$ 同色. 如果 d 也同色，done；否则看 2d；最后假设所有 d\textasciitilde{}kd 都和 x+id 不同色.

\notefigure{notes/250427-01.png}

这样一来，可以看成 d\textasciitilde{}kd(即 1\textasciitilde{}k) 被 c-1 染色. 如果 k 足够大，也一定存在 $x+6y=z$ 的同色解. 设使得任意 c-染色均存在 $x+6y=z$ 同色解的最小 $N=f(c)$. 这样一来，我们要求等差数列长度至少是 $6f(c-1)+1$, 因此 $f(c)\leq W(c,6f(c-1)+1)$. $\blacksquare$

推论：$x+2y+4z=w$ 是 Ramsey 方程

\notefigure{notes/250427-02.png}

\section{Rado Theorem}

$4x+7y+2z=2a+21b+3c$ 每个数可以拆成 $a=101^rq,101\not|\;q$, $\phi(a)=q\mod 101$ 即可(100染色)

假设有一个同色q解，两边约去等量的101因数，至少剩一个，假设剩余x,y,b，再模101，得到 $4q+7q=21q$, 得到 $q=0$，矛盾. （这里已经证了必要性）

\begin{quotation}
\textbf{Thm.} (Rado) 存在 $I\subseteq[n],J\subseteq[m]$, $\sum_Ia_i=\sum_Jb_j$, 当且仅当方程 $\sum_{i=1}^n a_ix_i=\sum_{j=1}^m b_jy_j$ 是 Ramsey 的.
\end{quotation}

充分性：每次假定两个变量相等，用其中一个替换另一个. 示例：$3x+2y+5z=a+6b+4c\to^{x=y,a=c}5x+5z=5a+6b\to^{z=b}5x=5a+b$, 如果 $5x=5a+b$ 有同色解，那么原方程也有同色解.

\notefigure{notes/250427-03.png}

\section{r-regular Ramsey}

$R(a_1,\dots,a_k;1)=\sum_{[k]}(a_i-1)+1$

我们估计 $R(15,26,31,17;5)$, 目标：给出的图对于任意5染色$\phi$存在15大小的同1色5子集，或者26大小的同2色5子集，...，或者存在17大小的同4色5子集。

假设挑出一个 $v$，剩下 $n-1$ 个点的4子集进行4染色，要求这个4染色是induced，i.e. $\sigma(\{a,b,c,d\})=i\iff\phi(\{v,a,b,c,d\})=i$，如果要求存在 $R(14,26,31,17;5)$ 个同1色，或...或 $R(15,26,31,16;5)$ 同4色, 那么对于这个4-子集同1色的完全图，或者同1色5子集至少14，或者同2色5子集至少26...如果加上v，由于4子集同1色完全图中，任意4-子集加上v就是同1色5子集，因此这个完全图加上v就是5子集同1色完全图。其他同理，满足条件。$R(15,26,31,17;5)\leq R(R(14,...;5),R(...,25,...;5),R(...,30,...;5),R(...,15;5);4)$

\section{Happy Ending Problem}

\begin{quotation}
\textbf{Thm.} (Esther Klein $\to$ Esther Szekeres) 平面上不三点共线的5个点必有4个点构成凸四边形.
\end{quotation}

考虑convex hull的形状

\begin{quotation}
\textbf{Thm.} (Erdos-Szekeres 1935) 任意k，存在足够大的n，平面上任意n个不三点共线的点，存在k点是convex的. 最小的n定义为 $N(k)$.
\end{quotation}

\textbf{Pf.} by Szekeres, 1931

\begin{quotation}
\textbf{Lem.} In $\mathbb{R}^2$, k pts (no 3 pts share a line) convex $\iff$ all 4 pts subset is convex.
\end{quotation}

用一个点对凸包其他顶点连线，进行三角形划分，每个点总在一个三角形内部.

因此考虑4超图的染色，四个点是凸的染y，非凸染b，$N\geq R(k,5;4)$ 即可. 意思是，对 $N$ 点的4超图进行2染色，必有k大小的同y色4子图，或者5大小的同b色4子图，但是我们知道5个点必有凸的四边形，因此第二种情况不可能，一定存在k大小的同y色4子图. $\blacksquare$

\textbf{Pf.} by Taski, 1971, 于某一次组合数学考试中憋出来

旋转使得没有两个点同x，左到右编号. 任意3个点 $i<j<k,a_i,a_j,a_k$ 有两种形态: 左转y/右转b，对这个3超图染yb色. $N\geq R(k,k;3)$, k个点中任意三个转向相同即可 $\blacksquare$

\textbf{Pf.} 1978

任意三个点连成三角形，内部奇数个点染y，偶数个点染b. $N\geq R(k,k;3)$, 意味着对这种染色，存在k个点，任意3子集都是y色或者任意3子集都是b色. 这k点中，不存在4点非凸子集，因为 $\text{奇}+\text{奇}+\text{奇}+1=\text{偶}$，$\text{偶}+\text{偶}+\text{偶}+1=\text{奇}$. 因此k点是凸的 (别忘了加上中间的点!).

\textbf{Pf.} by Erdos and Szekeres, 1935

\notefigure{notes/250427-04.png}

$N(p,q)$ 表示最小的 $N$, 平面上N个点(不共线，不共x)总存在p+2个上凸或者q+2个下凸. $N(0,\_)=2$, $N(n,1)=n+2$

$N_1+1=N(p-1,q),N_2+1=N(p,q-1)$. 考虑 $N_1+N_2+1$ 个点, 即证存在p+2个上凸或者q+2个下凸. 首先抓前 $N_1$ 个，如果存在q+2下凸，done; 如果存在p+1上凸，还差1; 如果运气够好可以加点，done;

(如果运气好，阴影部分存在点，就可以加上)

\notefigure{notes/250427-05.jpg}

如果没得加，我们反复对前$N_1+1$个点取存在p+2个上凸或者q+2个下凸。注意每次取完，把最后一个点从待取的list中去掉，如果满足前两个的条件就返回；假设取出来的序列都不满足条件，也就是说，都是p+1上凸，且每个都不一样（指存在不同的点, 但是可能有重复的点被用到），最终挑出来 $N_2+1$ 个这样的上凸链和（被去掉的）点. 如果运气够好，这 $N_2+1$ 个点中存在p+2上凸，done;

考虑这 $N_2+1$ 个点中存在q+1下凸，从左到右$v_1,\dots,v_{q+1}$, 回想 $v_1$ 是某一条上凸链的最后一个，而且这条上凸链运气不够好. 那么有q+2下凸，done.

\notefigure{notes/250427-06.jpg}

$N(p,q)-1\leq N_1+N_2=N(p-1,q)-1+N(p,q-1)-1\Rightarrow N(p,q)\leq{p+q\choose p,q}+1$

存在k个点上凸/下凸即可保证存在k个点凸。因此 $N(k)\leq N(k-2,k-2)\leq{2k-4\choose k-2}+1\sim 4^k$ (忽略这里的符号混用，你应该知道函数重载)

这个N(p,q)的上界还是紧的：对于 ${p+q\choose p,q}$ 个点，分裂成 ${p+q-1\choose p,q-1}+{p+q-1\choose p-1,q}=C_1+C_2$ 两个点集，其中 $C_1$ 没有p+2上凸也没有q+1下凸，$C_2$ 没有p+1上凸也没有q+2下凸（归纳假设），那么把这两个集合压扁，呈45度摆放足够远，$C_1$ 在下，$C_2$ 在上即可. 因此 $N(p,q)={p+q\choose p,q}+1$. 下面我们用这个思路证 $2^{k-2}+1\leq N(k)\leq N(k-2,k-2)$. 只需将2次幂化成组合数的和即可, 然后用上面的归纳.

\notefigure{notes/250427-07.png}

\notefigure{notes/250427-08.png}

\url{http://www.numdam.org/item?id=CM_1935__2__463_0}

\section{ramsey 数的下界}

$r(k,m)> (m-1)(k-1)$, 这是 Turan 定理的直接推论

考虑 $r(k,4)$, 随便 $p$-$1-p$ yb 染色, 按照 Erdos 的策略 $X_T$ for $|T|=k$, $Y_S$ for $|S|=4$, $Z=\sum X_T+\sum Y_S$.

\[
\mathbb{E}[Z]={n\choose k}(1-p)^{{k\choose 2}}+{n\choose 4}p^{6}
\]

我们可以期望这两项都小于1/2. 要使第二项满足，只需 $\frac{n^4}{4!}p^6<\frac{1}{2},\;\text{i.e. }p<\Big(\frac{12}{n^4}\Big)^{1/6}\sim n^{-2/3}$, 要使左边小于 1/2, 只需 ${n\choose k}e^{-p{k\choose 2}}<\frac{n^k}{k!}e^{-p(k-1)^2/2}<1/2$, 即只需 $k\ln n-k\ln k-\frac{p(k-1)^2}{2}+\ln 2<0$.

如果 $\epsilon>0$, 令 $n\sim k^{1.5-\epsilon}$, $p\sim n^{-2/3}\sim k^{-1+2\epsilon/3}$, 上式化为 $(0.5-\epsilon)k\ln k-k^{1+2\epsilon/3}<0$ 我们证明了 $n\sim k^{1.5-\epsilon}$ 存在没有坏集合的染色，即 $r(k,4)>k^{1.5-\epsilon}$
